% chapters/14/section-01.tex
\section{Jalan Pengolahan Aset}
\label{sec:pipeline-impor}

\subsection{Tiga Langkah dari Berkas menjadi Benda}

Perbedaan besar antara Unity dan blok Python terletak pada saat aset
dikerjakan. Pada Pygame, gambar dimuat ketika program dijalankan, sehingga
program harus disunting untuk mengubahnya. Pada Unity, aset diolah ketika
diimpor, dan hasilnya disimpan \cite{unity-manual-import}.

\begin{langkah}
  \item \textbf{Berkas disalin} ke dalam folder \code{Assets}. Unity memantau
        folder itu, sehingga berkas yang baru muncul langsung diolah.
  \item \textbf{Berkas diolah} menurut \istilah{importer setting} yang berlaku
        baginya, lalu hasilnya disimpan pada folder \code{Library}. Pengolahan
        itu hanya dilakukan sekali untuk setiap pengaturan, bukan setiap kali
        scene dijalankan \cite{unity-manual-import}.
  \item \textbf{Berkas dipakai} dengan menyeretnya dari \istilah{Project
        window} ke \istilah{Scene} atau ke komponen \istilah{Sprite Renderer}.
\end{langkah}

Akibat pemisahan itu penting untuk Anda ketahui: mengubah \istilah{importer
setting} berarti Unity mengolah ulang berkasnya, dan seluruh benda yang memakai
berkas itu \textbf{langsung berubah} tanpa perlu disentuh satu per satu. Itulah
sebabnya pengaturan impor menjadi tempat pertama yang perlu diperiksa ketika
sebuah aset tampak salah.

\subsection{Pengaturan Impor yang Paling Berpengaruh}

Pada aset dua dimensi, empat pengaturan berikut yang paling sering berpengaruh.

\begin{itemize}
  \item \textbf{\istilah{Sprite Mode}.} Menentukan apakah berkas gambar berisi
        satu gambar atau banyak gambar yang tersusun sebagai lembar. Nilai
        \istilah{Single} untuk gambar tunggal, dan \istilah{Multiple} untuk
        \istilah{sprite sheet}.
  \item \textbf{\istilah{Pixels Per Unit}.} Menentukan berapa piksel gambar
        yang setara dengan satu satuan pada scene. Nilai inilah yang mengaitkan
        ukuran gambar dengan ukuran level yang telah Anda tetapkan pada
        Bagian~13.1.
  \item \textbf{\istilah{Filter Mode}.} Menentukan cara gambar dihaluskan ketika
        diperbesar atau diperkecil. Nilai \istilah{Point} mempertahankan
        ketajaman piksel, dan biasanya dipilih untuk gambar bergaya piksel.
  \item \textbf{\istilah{Compression}.} Menentukan cara gambar dipadatkan saat
        program dijalankan. Nilai \istilah{None} menjaga mutu gambar, sedangkan
        nilai lain menghemat ruang dengan mengorbankan sedikit ketajaman.
\end{itemize}

\begin{peringatanbox}
Jangan mengubah ukuran berkas gambar di luar Unity, misalnya dengan program
penyunting gambar. Apabila berkasnya berubah, Unity akan mengolahnya ulang, dan
hasilnya untuk seluruh benda yang memakainya berubah pula. Ukuran gambar
sebaiknya diubah melalui pengaturan \istilah{Pixels Per Unit} pada Unity, agar
perubahannya terkendali dan dapat dikembalikan.
\end{peringatanbox}

\subsection{\istilah{Sprite Sheet} dan Pemotongannya}

Sebuah \istilah{sprite sheet} memuat beberapa gambar dalam satu berkas, dan
pemotongannya diatur melalui \istilah{Sprite Editor}. Cara itu dipakai untuk
gambar yang berubah-ubah, misalnya gerak berjalan yang dipecah menjadi beberapa
bingkai.

Kebiasaan yang menghemat waktu: tentukan ukuran setiap bingkai secara
seragam, misalnya 64 kali 64 piksel, sehingga pemotongannya dapat dilakukan
secara otomatis menurut ukuran itu. Ukuran yang tidak seragam menuntut
pemotongan satu per satu, dan itu melelahkan serta mudah salah.
